\quellseite{23}
\chapter{Introduction}

Henry I. Schmidt, \emph{History of Education}, Part i: \emph{History of Education, Ancient and Modern}; Part ii: \emph{A plan of Culture and Instruction based on Christian Principles, and designed to aid in the right Education of youth, Physically, Intellectually, and Morally}, 1842. L. Kraussold, \emph{Katechetik fuer Schule und Kirche}, 1843, ²1880. Chr. Palmer, \emph{Evangelische Katechetik}, 1844, ⁶1881. *G. v. Zezschwitz, \emph{System der christlich-kirchlichen Katechetik}: vol. i: \emph{Der Katechumenat oder die Lehre von der kirchlichen Erziehung}, 1863; vol. ii: \emph{Die Lehre vom kirchlichen Unterricht nach Stoff und Methode}: part i: \emph{Der Katechismus oder der kirchliche Unterrichtsstoff}, 1864, ²1872; part ii: \emph{Die Katechese oder die kirchliche Unterrichtsmethode}; div. i: \emph{Der akroamatisch-positive Bibelunterricht}, 1869, ²1874; div. ii: \emph{Die erotematische Unterrichtsform}, 1872. Ziegler, \emph{Catechetics, Historical, Theoretical, Practical}, 1873. F. W. Schuetze, \emph{Praktische Katechetik fuer evangelische Seminare und Lehrer}, 1876, ²1883. R. Kuebel, \emph{Katechetik}, 1877. Theodor Harnack, \emph{Katechetik und Erklaerung des Kleinen Katechismus}, 1882. G. v. Scheele, \emph{Kirchliche Katechetik}, 1886. *K. Buchrucker, \emph{Grundlinien der kirchlichen Katechetik}, 1889. E. Sachsse, \emph{Die Lehre von der kirchlichen Erziehung nach evangelischen Grundsaetzen}, 1897. *Fr. Zange, \emph{Evangelischer Religionsunterricht}, 1897, ²1913. A. Eckert, \emph{Der erziehende Religionsunterricht in Schule und Kirche}, 1899, ²1915. M. v. Nathusius, \emph{Handbuch des kirchlichen Unterrichts}, 1903. O. Baumgarten, \emph{Neue Bahnen. Der Unterricht in der christlichen Religion im Geist der modernen Theologie}, 1903. E. Thraendorf, \emph{Allgemeine Methodik des Religionsunterrichts}, 1903, ⁵1912. *G. A. Coe, \emph{Education in Religion and Morals}, 1904. J. Gottschick, \emph{Homiletik und Katechetik}, 1908. Joh. Berndt, \emph{Methodik des Unterrichts in der evangelischen Religion}, 1909. J. Smend, \emph{Der evangelische Religionsunterricht auf hoeheren Schulen}, 1910. *R. Kabisch, \emph{Wie lehrt man Religion?} 1910, ²1912. G. H. Gerberding, \emph{The Lutheran Catechist}, 1910. J. H. Herzer, \emph{Ev.-Luth. Katechetik}, 1911. C. W. Hertzler, \emph{Die religioes-sittliche Erziehung der kirchlichen Jugend}, 1911. K. Knoke, \emph{Recht und Pflicht der evangelischen Kirche hinsichtlich der religioesen Unterweisung ihrer heranwachsenden Jugend}, 1912. Chr. Buerckstuemmer, \emph{Der Religionsunterricht in der Volksschule}, 1913.
\quellseite{24}
J. Steinbeck, \emph{Lehrbuch der christlichen Jugenderziehung}, 1914. *G. Hodges, \emph{The Training of Children in Religion}, 1911, 1917. *E. Hershey Sneath, George Hodges, and Henry H. Tweedy, \emph{Religious Training in the School and Home}, 1917. *G. H. Betts, \emph{How To Teach Religion}, 1919. *E. Pfennigsdorf, \emph{Wie lehren wir Evangelium?} 1921, ²1925. (S. M. Cavert), \emph{The Teaching Work of the Church}, 1923. P. E. Kretzmann, \emph{The Teaching of Religion}, 1925. Compare also the chapters on Catechetics by G. v. Zezschwitz in: Strack-Zoeckler, \emph{Handbuch der theologischen Wissenschaften}, part iv, 1882, ³1890. K. Knoke, \emph{Grundriss der praktischen Theologie}, 1886, ⁴1896. A. Krauss, \emph{Lehrbuch der praktischen Theologie}, 1890, 1893. *E. Chr. Achelis, \emph{Lehrbuch der praktischen Theologie}, 1890, ³1911. E. Chr. Achelis, \emph{Grundriss der praktischen Theologie}, 1893, ⁵1903. F. Niebergall, Praktische Theologie, 1918–19. M. Schian, \emph{Grundriss d. Prakt. Theologie}, 1923. Of textbooks on Pedagogics we mention: Chr. Palmer, \emph{Evangelische Paedagogik}, 1853, ⁴1869. T. Ziller, \emph{Vorlesungen ueber allgemeine Paedagogik}, 1876, ³1892. G. v. Zezschwitz, \emph{Lehrbuch der Paedagogik}, 1882. *W. Ostermann-L. Wegner, \emph{Lehrbuch der Paedagogik}, 1882, ⁵1910. J. Chr. G. Schumann-G. Voigt, \emph{Lehrbuch der Paedagogik}, about 1882, ¹²1904. H. Spencer, \emph{Education, intellectual, moral, and physical}, 1883. Jos. Payne, \emph{Lectures on the Science and Art of Education}, 1884. E. Thring, \emph{Theory and Practice of Teaching}, 1885. J. G. Fitch, \emph{Lectures on Teaching}, 1886. J. Johonnot, \emph{Principles and Practice of Teaching}, 1886. *Ch. de Garmo, \emph{The Essentials of Method}, 1889. D. P. Page, \emph{Theory and Practice of Teaching}, 1893. *W. James, \emph{Talks to Teachers on Psychology}, 1900. K. Knoke, \emph{Grundriss der Paedagogik}, 1892, 1902. *K. Heilmann, \emph{Handbuch der Paedagogik}, 1897, 1908. W. Rein, \emph{Paedagogik im Grundriss}, 1900, ⁴1908; English by C. C. and Ida J. Van Liew, 1895. W. Rein, \emph{Paedagogik in systematischer Darstellung}, 1902–1906. P. Barth, \emph{Die Elemente der Erziehungs- und Unterrichtslehre}, 1906, ³1911. F. Paulsen, \emph{Paedagogik}, 1911. K. A. Schmid, \emph{Geschichte der Erziehung}, 5 parts, 11 volumes, 1884 ff. K. A. Schmid, \emph{Enzyklopaedie des gesamten Erziehungs- und Unterrichtswesens}, 2. ed., 8 volumes. W. Rein, \emph{Enzyklopaedisches Handbuch der Paedagogik}, 2. ed., 10 volumes. \emph{Sonnenschein’s Cyclopedia of Education}, ed. by Fletcher, 1892. \emph{Monroe’s Cyclopedia of Education}, 1911–1914. W. S. Monroe, \emph{Bibliography of Education}, 1897, 1907. P. Monroe, \emph{Textbook in the History of Education}, 1905. *F. P. Graves, \emph{History of Education before the Middle Ages}, 1909; \emph{History of Education during the Middle Ages and the Transition to Modern Times}, 1909–1910.

\quellseite{25}
\section{§ 1. The Definition of Catechetics}

G. v. Zezschwitz i, pp. 10–34; ii, 1, pp. 27–43. Achelis ii, pp. 278–281; 383–386. Sachsse, pp. 301–302.

The term Catechetics (κατηχητικὴ τέχνη) is derived from the compound verb κατηχέω. This, in turn, is developed from the simple verb ἠχεῖν = to sound, ring, peal, not from the noun ἠχώ = echo, and thus signifies a descending sound, not a reverberating sound, an echo, or a response, for which the Greeks employed ἀντηχεῖν. The etymology of the word Catechetics therefore disproves the explanation widely spread by Melanchthon though not originated by him, according to which this term indicates instruction in the form of questions and answers.\footnote{Compare M. Reu, \emph{Quellen zur Geschichte des kirchlichen Unterrichts}, i, 2, p. 17; F. Cohrs, \emph{Evangelische Katechismusversuche vor Luthers Enchiridion}, iii, p. 20.} Originally intransitive, the verb subsequently received transitive force: κατηχεῖσθαι = receive information, κατηχεῖν τινα = to reach one by sound from above (as from a desk), to impart knowledge by the transmission of sound from lip to ear; frequently it had the connotation of the superficial, inchoate, elementary, rudimentary. In the general sense of imparting oral information Luke uses the verb in his Gospel (1: 4) and in Acts (18: 25; 21: 21); Paul uses it in the more definite sense of giving oral instruction in religion (Rom. 2: 18; 1 Cor. 14: 19; Gal. 6: 6). Inasmuch as oral instruction in the fundamentals of Christianity was the principal feature of the preparation for baptism, κατηχέω was used as the technical term for such preparation and subsequently came to signify all fundamental religious instruction imparted by the Church. Catechetics, accordingly, is an exposition of the principles of religious instruction and a practical method of equipping the catechist.

\quellseite{26}
\section[§ 2. The Teaching Function of the Church]{§ 2. The Teaching Function of the Church and the Necessity of Catechetics}

*Th. Harnack, pp. 1–8. Achelis ii, pp. 281–282. Buchrucker, pp. 64–65. *Kabisch, pp. 1–102. \emph{Religionsunterricht? 80 Gutachten. Ergebnis einer von der Vereinigung fuer Schulreform in Bremen veranstalteten allgemeinen deutschen Umfrage}, 1905. W. Rein, \emph{Stimmen zur Reform des Religionsunterrichts}, 1904. ²1906. G. A. Coe, \emph{Education in Religion and Morals}, pp. 21–32. Sneath, Hodges, and Tweedy, \emph{Religious Training in School and Home}, pp. 1–33.

Religious instruction is one of the indispensable obligations of the Church. The very nature and the task of the church suggest the necessity of imparting religious instruction. The Church is to spread and grow. Through her missionary activity she endeavors to arouse those without to a longing for salvation and to desire membership; but her nature as the communion of believers prevents her from receiving into fellowship any but those who share her faith and confession. Thus it becomes evident that the Church must carefully instruct all applicants in her faith and confession. The need of such instruction results, moreover, from baptism as the divinely ordained Sacrament of initiation into the Church. Baptism necessitates instruction concerning its character and import; circumstances determine when this instruction will be given, whether before or after the Sacrament, but it must be given. No adult whether Pagan, Jew, or antibaptist sectarian, in whom the missionary has awakened the wish to join the Church, will be admitted to baptism without previous instruction concerning baptism and the faith; and the baptism of infants dare not take place at all if there is no guarantee whatsoever of subsequent instruction. Finally, religious instruction is explicitly required in Holy Writ. The missionary command involves the training
\quellseite{27}
for a Christian life of those who have been initiated into the Church by baptism (Matt. 28:18ff.); even more, Jesus impressed upon Peter not only the guarding and feeding of the sheep (\emph{πρόβατα}, \emph{προβάτια}), but also the feeding of the lambs (\emph{ἀρνία}) belonging to the fold (John 21:15ff.), a process to which instruction belongs as an integral part. Since Rousseau’s times († 1778) the demand has indeed frequently been voiced that children should be given no religious information before the age of 16 “in order not to forestall an independent decision”; what is actually demanded under this pretext, however, is an irreligious education inasmuch as innate evil is given an opportunity for unimpeded development; besides, the parents and the Church are expected to withhold their own most precious treasure from their children, and this at a time when in every respect the foundations are laid for the future.

E. H. Sneath, George Hodges and H. H. Tweedy write in \emph{Religious Training in the School and Home}:
\begin{quote}
“To postpone the work of religious education until the preparatory school and college is as rash as it is foolish. ‘Let a child wait until he is grown and then choose his own religion’, said an English statesman in the hearing of Coleridge. Coleridge made no reply, but led the speaker out into his garden. Looking around upon the bare ground he said quietly: ‘I have decided not to put out any flowers and vegetables this year, but to wait till August and let the garden decide for itself whether it prefers weeds or strawberries!’ The blind optimism which seems to feel that the moral and religious training of the child will care for itself, ends all too soon in disaster. The business of growing good men and women will no more care for itself than the business of making a fortune or winning a success in any profession will care for itself” (p. 19).
\end{quote}

Nor is it the mere principle of instruction as such that is in question, but rather the purposeful and systematic instruction and training. Instruction bearing a general Christian impress might be conceivable as the result of Christian association and a life within the pale of the Church. But
\quellseite{28}
quite aside from the failure of many Christians to possess an adequate understanding of the Way of Salvation or to give proper ethical expression to it, such incidental instruction would leave a complete and connected teaching of essential truth a matter of chance. Gaps, misunderstandings, aberrations, would be all but inevitable while purposeful influence and coherent effort would be absent altogether. It, therefore, is the imperative obligation of the Church, as an organization, to pay most earnest attention to this work—her very future is at stake—, to prescribe an orderly process of instruction, to gain greater clearness concerning the aim, the matter, and the method, of such instruction, and most conscientiously to prepare her future servants, in both theory and practice, for a successful discharge of their teaching duties. The history of the Church is witness of the wholesome effect ever produced upon her inner development through the sacred performance of this sacred duty, and of the baneful effects wherever it has been neglected.

\quellseite{29}
\section{§ 3. The Division of Catechetics}

Inasmuch as the Church has had to give spiritual instruction from the time of her origin, the first element of the catechetical discipline will be a survey of its \emph{history}; thereby the experience of the past is made available for the present, and a clearer and ampler appreciation and discharge of her teaching function is rendered possible. In order to discover the proper manner of training, the \emph{pupil’s mental faculties and their gradual development} must be studied. Furthermore, it will be necessary definitely to formulate the \emph{aim} of instruction, to decide what \emph{materials} must be employed in order to reach this \emph{aim} and how they are to be \emph{distributed} over the accessible educational agencies, and to study the \emph{method} which best assures the attainment of the \emph{aim}. And, finally, we must inquire when and how such instruction is to be \emph{concluded}.
